% Generated draft: requires mathematical and source review.
\documentclass[11pt]{article}
\usepackage[margin=25mm]{geometry}
\usepackage{fontspec}
\setmainfont{Latin Modern Roman}
\usepackage{amsmath,amssymb,mathtools,unicode-math}
\setmathfont{Latin Modern Math}
\usepackage{enumitem,needspace,xurl,hyperref}
\setlength{\parindent}{0pt}
\setlength{\parskip}{5pt}
\setlength{\emergencystretch}{3em}
\newcommand{\N}{\mathbb{N}}
\newcommand{\Z}{\mathbb{Z}}
\newcommand{\Q}{\mathbb{Q}}
\newcommand{\R}{\mathbb{R}}
\newcommand{\C}{\mathbb{C}}
\newcommand{\F}{\mathbb{F}}
\newcommand{\A}{\mathbb{A}}
\newcommand{\PP}{\mathbb{P}}
\newcommand{\E}{\mathbb{E}}
\newcommand{\eps}{\varepsilon}
\newcommand{\dd}{\,\mathrm{d}}
\begin{document}
\section*{UnsolvedMath ID: 700847}
\textbf{Source problem code:} catalogue entry
\subsection{Definitions and mathematical statement}
Let \(G\) be a group with identity \(1_G\). A function \(f\colon G\to\C\) is called \emph{normed} if \(f(1_G)=1\), \emph{central} if
\[
 f(gh)=f(hg)\qquad(g,h\in G),
\]
and \emph{positive-definite} if, for every integer \(n\geq1\), every \(g_1,\ldots,g_n\in G\), and every \(c_1,\ldots,c_n\in\C\),
\[
 \sum_{k=1}^{n}\sum_{l=1}^{n} f(g_k^{-1}g_l)\overline{c_k}c_l\in\R_{\geq0}.
\]
A group is \emph{locally finite} if each subgroup generated by finitely many elements is finite, and \emph{simple} if it is nontrivial and its only normal subgroups are \(\{1_G\}\) and \(G\).

The supplied classification question is: classify the infinite simple locally finite groups \(G\) admitting a normed, central, positive-definite function \(f\colon G\to\C\).

\textbf{Unresolved qualification.} These three conditions alone do not express the intended distinction: on every group, both the constant function \(f(g)=1\) and the function
\[
 \delta_{1_G}(g)=\begin{cases}1,&g=1_G,\\0,&g\neq1_G\end{cases}
\]
satisfy them. The supplied source nevertheless asserts that simple Chevalley groups have no such functions, whereas locally matrix (or stable) classical groups over finite fields possess such functions. These are supplied source assertions, and the nonexistence assertion is incompatible with the displayed conditions alone. An additional restriction on the admissible functions appears to be missing; the record does not specify it. Consequently, a precise formulation of the intended classification remains unresolved, and no additional restriction is imposed here.

The source describes the question as motivated by the theory of \(C^*\)-algebras and identifies the references below as motivation and partial results. Its status label is \emph{open}; this is supplied registry metadata, not an independently verified assessment.
\subsection{Short English statement}
The source asks which infinite simple locally finite groups admit normalized, central, positive-definite complex functions. Its stated nonexistence claim conflicts with the given definition, since the constant function and the identity indicator always qualify. A missing restriction must be recovered before the intended classification can be stated precisely.
\subsection{Sources}
\begin{itemize}
\item Local UnsolvedMath registry, supplied complete source record, permanent ID \(700847\), exact source problem code \texttt{catalogue entry}, title \emph{Kourovka Notebook Problem 12.28}. The record was created and updated on \texttt{2026-10-16T14:00:00Z} and labels the problem open; this label is supplied metadata, not independently verified.
\item \emph{The Kourovka Notebook (Unsolved Problems in Group Theory)}, problem 12.28, 12th issue (1992); supplied edition reference: 21st issue edition (2026), \texttt{arXiv:1401.0300v48}. The supplied collection description names E. I. Khukhro and V. D. Mazurov as editors. The proposer is supplied as A. E. Zalesski with incomplete trailing accent markup; the proposed year is 1992.
\item A. M. Vershik and S. V. Kerov, \emph{J. Sov. Math.}, \textbf{38} (1987), 1701--1733, \(\S\,9\). Cited by the supplied source for motivation from \(C^*\)-algebras.
\item F. Leinen and O. Puglisi, \emph{J. Pure Appl. Algebra}, \textbf{208} (2007), 1003--1021; and \emph{J. London Math. Soc.}, \textbf{70} (2004), 678--690. Both are identified by the supplied source as containing partial results.
\end{itemize}
\end{document}
